% META: {"id": "TCOMPL-BAYES-ME-005", "chapitre": "TCOMPL-INFERENCE-BAYESIENNE", "type_objet": "methode", "capacites_codes": ["C5"], "methodes": ["M5"], "mode_creation": "ex_nihilo", "sources_inspiration": [], "status": "needs_review"}
% BEGIN-VERIFY
% from sympy import Rational, symbols, simplify, diff
% x = symbols('x', positive=True)
% Se, Sp = Rational(9,10), Rational(19,20)
% V = Se*x/(Se*x + (1-Sp)*(1-x))
% f = 18*x/(17*x+1)
% assert simplify(V - f) == 0
% assert f.subs(x, Rational(1,100)) == Rational(2,13)
% assert f.subs(x, Rational(1,2)) == Rational(18,19)
% assert simplify(diff(f, x) - 18/(17*x+1)**2) == 0
% END-VERIFY
\begin{fichemethode}{M5}{Etudier la VPP en fonction de la prevalence}

\quandUtiliser{%
  Utiliser cette methode lorsque l'enonce demande d'exprimer la valeur
  predictive positive comme FONCTION de la prevalence $p$, d'etudier ses
  variations, et d'interpreter (depistage de masse ou population ciblee).
}

\pasApas{%
  \begin{enumerate}
    \item Poser $V(p) = \dfrac{\text{Se} \cdot p}{\text{Se} \cdot p +
      (1-\text{Sp})(1-p)}$ pour $p \in\ ]0\,;\,1[$, ou Se et Sp sont fixees.
    \item Simplifier l'expression en une fonction homographique de $p$.
    \item Etudier les variations (derivee, ou croissance d'une fonction
      homographique de determinant positif) : $V$ est croissante.
    \item Calculer $V$ pour les valeurs de $p$ demandees (fraction exacte puis
      arrondi) et noter les limites : $V(p) \to 0$ quand $p \to 0$,
      $V(p) \to 1$ quand $p \to 1$.
    \item Interpreter : plus la maladie est rare, moins un test positif est
      probant — d'ou l'interet de cibler une population a risque avant de
      tester.
  \end{enumerate}
}

\exempleRedige{%
  Avec $\text{Se} = 0{,}9$ et $\text{Sp} = 0{,}95$, exprimer la VPP en fonction
  de la prevalence $p$, puis comparer $p = 1\,\%$ et $p = 50\,\%$.

  \medskip
  \commentaireMarge{Multiplier numerateur et denominateur par 20.}%
  $V(p) = \dfrac{0{,}9\,p}{0{,}9\,p + 0{,}05\,(1-p)} = \dfrac{18p}{17p + 1}$.

  \medskip
  $V'(p) = \dfrac{18}{(17p+1)^2} > 0$ : $V$ est croissante sur $]0\,;\,1[$.

  \medskip
  \commentaireMarge{La meme qualite de test, deux fiabilites tres differentes.}%
  $V(0{,}01) = \dfrac{2}{13} \approx 0{,}15$ et
  $V(0{,}5) = \dfrac{18}{19} \approx 0{,}95$ : en population generale (maladie
  rare), un positif est peu probant ; dans une population ciblee ou la
  prevalence atteint 50\,\%, le meme test positif devient tres fiable.
}

\pieges{%
  \begin{itemize}
    \item \textbf{Piege 1.} Croire que la VPP ne depend que de la qualite du
      test : a Se et Sp fixees, elle varie fortement avec la prevalence.
    \item \textbf{Piege 2.} Oublier le facteur $(1-p)$ devant $(1-\text{Sp})$
      dans le denominateur.
    \item \textbf{Piege 3.} Conclure « le test est mauvais » quand la VPP est
      faible : c'est la rarete de la maladie, pas le test, qui est en cause.
  \end{itemize}
}

\verifier{%
  Controler $0 < V(p) < 1$, retrouver une valeur particuliere par l'arbre
  pondere complet (fiche M4), et verifier les comportements limites
  ($V \to 0$ en $0$, $V \to 1$ en $1$).
}

\sEntrainer{\refExos{M5}}

\end{fichemethode}
